\documentclass[11pt]{article}
\usepackage[utf8]{inputenc}
\usepackage[T1]{fontenc}
\usepackage[spanish,es-noshorthands]{babel}
\usepackage[margin=2.5cm]{geometry}
\usepackage{amsmath,amssymb}
\usepackage{graphicx}
\usepackage{booktabs}
\usepackage{xcolor}
\usepackage[colorlinks=true,linkcolor=blue,citecolor=blue]{hyperref}
\usepackage{caption}
\usepackage{parskip}
\usepackage[expansion=false]{microtype}
\sloppy

\title{\textbf{¿Sigue funcionando \texttt{search\_best\_beta\_mejorado} fuera de su rango de conductancia calibrado?}\\
\large Validación de la búsqueda de beta y de la fórmula analítica para ratios $G_{\max}/G_{\min}\in\{2,20,100\}$}
\author{Walter Quiñonez}
\date{Reporte generado con Claude Code \textperiodcentered\ 10 de septiembre de 2026}

\begin{document}
\maketitle

\begin{quote}\itshape
\texttt{MemDNNMejorado.search\_best\_beta()} (\texttt{mejora\_search\_best\_beta/}) y la fórmula
analítica que centra su grilla (\texttt{beta\_pred\_formula}) se calibraron y validaron con el
rango de conductancia fijo de todo el proyecto ($R_{\text{low}}=1000\,\Omega$,
$R_{\text{high}}=10000\,\Omega \Rightarrow$ ratio $G_{\max}/G_{\min}=10$). Esta corrida prueba
si siguen funcionando cuando ese rango cambia (ratio $\in\{2,20,100\}$, con INL de la curva
P/D casi constante para aislar el efecto del rango de conductancia). Conclusión: \textbf{sí}
--- el regret del beta elegido se mantiene por debajo de 0.52pp en los tres ratios, incluso en
ratio=2, donde la fórmula subestima el beta óptimo real por un factor $\sim$5.9$\times$ y la
red de seguridad (expansión automática de la grilla) es la que sostiene la precisión.
\end{quote}

\tableofcontents

\section{Pregunta y metodología}

Para cada ratio $G_{\max}/G_{\min}\in\{2,20,100\}$ ($R_{\text{low}}=1000\,\Omega$ fijo,
$R_{\text{high}}=R_{\text{low}}\cdot\text{ratio}$; ratio=10 ya está validado en
\texttt{mejora\_search\_best\_beta/}, no se repite acá):

\begin{enumerate}
\item \textbf{Referencia} (independiente de la fórmula): grilla ancha y fija de 10 betas
log-espaciados en $[20,8000]$, evaluada con el protocolo mínimo validado en
\texttt{analisis\_epochs\_series\_minimos/} (Fase 3: series=8, k\_folds=3, epochs=8). El beta
con mayor accuracy media de esa grilla es el ``óptimo real'' de referencia para ese ratio.
\item \textbf{Búsqueda mejorada}: se corre \texttt{search\_best\_beta} 3 veces con semillas
distintas, y se compara el beta elegido contra la referencia (regret = accuracy de la
referencia en su óptimo $-$ accuracy de la referencia en el punto de grilla más cercano al beta
elegido).
\end{enumerate}

INL$_{\text{pot}}$ se mantuvo prácticamente constante entre ratios (0.0101 para ratio=20 y
ratio=100; documentado entre 0.0065 y 0.0101 para ratio 2..100 en el script), para aislar el
efecto del rango de conductancia del efecto de no linealidad.

Código: \texttt{estudiar\_rangos\_conductancia.py} (no modifica ningún archivo original del
proyecto). Datos completos, verificables: \texttt{resultados\_rangos\_conductancia.json} y
\texttt{referencia\_ratio\{2,20,100\}.json} en esta misma carpeta.

\section{Resultados}

\subsection{Grillas de referencia}

\begin{table}[htbp]
\centering
\begin{tabular}{@{}lrr@{}}
\toprule
Ratio $G_{\max}/G_{\min}$ & $\beta_{\text{óptimo referencia}}$ & acc.\ óptima referencia \\
\midrule
2   & 557.98 & 0.9074 \\
20  & 286.75 & 0.9078 \\
100 & 286.75 & 0.9079 \\
\bottomrule
\end{tabular}
\caption{Óptimo de la grilla ancha y fija de 10 betas (log-espaciados en [20,8000]), protocolo
mínimo. Ratio 20 y 100 caen en el mismo punto de esa grilla (resolución de la grilla, no
necesariamente el mismo óptimo continuo); el accuracy en ese punto es prácticamente idéntico
entre ambos.}
\end{table}

El óptimo real se corre hacia betas más chicos a medida que aumenta el ratio (más rango de
conductancia disponible por bit de resolución), consistente con lo esperado.

\subsection{Búsqueda mejorada: predicción de la fórmula vs.\ beta elegido}

\begin{table}[htbp]
\centering
\begin{tabular}{@{}lrrrrr@{}}
\toprule
Ratio & $\beta_{\text{pred (fórmula)}}$ & factor vs.\ real & expansiones & regret medio & regret peor caso \\
\midrule
2   & 3294.9 & 5.91$\times$ & 1 (las 3 corridas) & 0.09 pp & 0.27 pp \\
20  & 315.6  & 1.10$\times$ & 0                   & 0.15 pp & 0.23 pp \\
100 & 271.4  & 0.95$\times$ & 0                   & 0.17 pp & 0.51 pp \\
\bottomrule
\end{tabular}
\caption{3 repeticiones (semillas nuevas) por ratio. ``factor vs.\ real'' =
$\beta_{\text{pred (fórmula)}}/\beta_{\text{óptimo referencia}}$. Detalle de las 9 corridas en
\texttt{resultados\_rangos\_conductancia.json}.}
\end{table}

\textbf{Ratio=2 es el caso interesante}: la fórmula (ajustada y validada solo en ratio=10)
subestima el beta óptimo por casi 6$\times$ --- la extrapolación fuera de su rango de
calibración falla claramente acá, mucho más que en ratio=20 o 100 (factores 1.10$\times$ y
0.95$\times$, dentro del error ya documentado de la fórmula en \texttt{formula\_beta\_optimo\_propuesta/}).
Sin embargo, la red de seguridad de \texttt{search\_best\_beta\_mejorado} (expansión automática
del margen cuando el óptimo empírico cae en el borde de la grilla angosta) se activó en las 3
corridas de ratio=2, y el beta finalmente elegido (653.8 u 823.7) quedó igual de cerca del
óptimo real (557.98) que en los ratios donde la fórmula no falló: regret peor caso 0.27pp,
comparable a los otros ratios. El costo de esa expansión es evaluar el doble de puntos de
grilla: las búsquedas de ratio=2 tardaron $\approx$442s contra $\approx$235--248s de ratio=20/100
sin expansión.

\subsection{En contexto: los 4 ratios probados hasta hoy}

\begin{table}[htbp]
\centering
\begin{tabular}{@{}lrr@{}}
\toprule
Ratio $G_{\max}/G_{\min}$ & regret peor caso & fuente \\
\midrule
2 (este reporte)             & 0.27 pp & \texttt{resultados\_rangos\_conductancia.json} \\
10 (calibración original)    & 0.45 pp & \texttt{mejora\_search\_best\_beta/resultados\_validacion.json} \\
20 (este reporte)            & 0.23 pp & \texttt{resultados\_rangos\_conductancia.json} \\
100 (este reporte)           & 0.51 pp & \texttt{resultados\_rangos\_conductancia.json} \\
\bottomrule
\end{tabular}
\caption{El dato de ratio=10 viene de un protocolo de referencia distinto (ground truth de 100
épocas/20 series/5 folds sobre 2 combos INL, no la grilla mínima de 10 puntos usada acá) -- no
es estrictamente comparable número a número, pero da la misma conclusión cualitativa: regret
peor caso siempre por debajo del umbral de 1pp que \texttt{analisis\_epochs\_series\_minimos/}
estableció como aceptable.}
\end{table}

\section{Conclusión}

La búsqueda mejorada de beta \textbf{sigue funcionando bien fuera del rango de conductancia con
el que se calibró} (ratio=10), en los tres ratios adicionales probados (2, 20, 100): el regret
del beta elegido contra la referencia nunca superó 0.51pp, muy por debajo del umbral de 1pp ya
establecido como aceptable en este proyecto. Esto se sostiene incluso cuando la fórmula
analítica que centra la grilla de búsqueda extrapola mal (ratio=2, factor $\sim$5.9$\times$):
el mecanismo de expansión automática del margen compensa el error de la fórmula al costo de
duplicar el número de puntos de grilla evaluados (y por lo tanto el tiempo de búsqueda) en ese
régimen.

\section{Limitaciones}

\begin{itemize}
\item La referencia usa una grilla fija de solo 10 betas log-espaciados (protocolo mínimo, no
el ground truth completo de \texttt{data\_beta\_grid\_search\_SP/}), igual que se documenta en
\texttt{estudiar\_rangos\_conductancia.py}; es suficiente para el protocolo ya validado pero no
un ground truth perfecto.
\item Ratio 20 y 100 dieron el mismo punto óptimo de referencia (286.75) -- posible efecto de
la resolución de la grilla de referencia (10 puntos entre 20 y 8000) más que de que el óptimo
continuo sea idéntico.
\item El dato de ratio=10 usado como contexto viene de un protocolo distinto (ver tabla de la
sección anterior), no es una corrida nueva con el mismo protocolo que ratio=2/20/100.
\item Solo 3 repeticiones por ratio (igual que se especificó en el script), no las 5 usadas en
la validación original de \texttt{mejora\_search\_best\_beta/}.
\item Sigue sin optimizarse un beta distinto por capa (mismo alcance que el resto del proyecto).
\end{itemize}

\end{document}
